\documentclass[11pt]{article}
\usepackage{ap-lang}
\title{Part 3: Reasoning and Organization: Question Set}
\author{Paul Wang}
\date{November 2026}
\docpart{Part 3 \textperiodcentered\ Question Set}

\begin{document}
\maketitle

\begin{abstract}
Three exam-style passages with multiple-choice questions, followed by a full synthesis essay (Question 1) with six sources, two of them quantitative. Passage 1 is an essay that moves from the story of a stopped clock to an argument about maintenance; Passage 2 is a lecture to legislators defending basic research. Both turn on structure, so most questions ask how a paragraph advances the line of reasoning or how the passage is developed. Passage 3 is a writing set on a student's local history essay that needs reordering. The passages, sources, writers, and places are fictional and were written for this practice set. Allow about 35 minutes for the multiple choice and 55 minutes for the essay (15 to read the sources, 40 to write).
\end{abstract}

\scoringline{MCQ: \blank\ / \qtotal \qquad Essay: \blank\ / 6 \qquad Target: 19+ MCQ and 4+ essay \qquad Pace: 80 s per MCQ}

\section{Passage 1: The courthouse clock (reading)}

\begin{passage}{Passage 1}{The following passage is from an essay published in 2023 in a regional newspaper's weekend magazine.}
\pn{1} The clock on the Brennan County courthouse stopped at 4:17 on a Tuesday in March of 1998, and for eleven years it told the correct time twice a day. People made jokes about it. Visitors photographed it. A local bakery sold a cookie shaped like its face, the hands frozen at 4:17, and the cookie outsold everything else in the case.

\pn{2} When the county finally repaired the clock, the job took one clockmaker four days and cost less than the annual budget for the courthouse's holiday lights. The problem had been a worn brass gear about the size of a quarter. A replacement had to be cut by hand, because the company that made the original had closed in 1931.

\pn{3} I tell this story because it is a small version of a large habit. We are eager to build and reluctant to maintain. A new bridge gets a ribbon-cutting, a speech, and a photograph in the paper; repainting the old one gets a detour sign. Politicians who would never miss the opening of a library will not attend the replacement of its roof. Nobody names a boiler after a donor.

\pn{4} The habit is expensive. Engineers have a rough rule that every dollar of maintenance deferred becomes several dollars of repair later, as a leak becomes rot and rot becomes a collapsed ceiling. A pothole patched in April costs a few dollars of asphalt; the same pothole in November, widened by a season of freezing and thawing, costs a repaving crew.

\pn{5} Maintenance is invisible precisely when it succeeds. A well-kept bridge looks exactly like a bridge, and so the work that keeps it standing is never seen and never praised. That invisibility is the real problem, and it suggests the real solution: make the work visible. Publish the maintenance schedule the way we publish ribbon-cuttings. Put the names of the crews who reseal the roof on the plaque beside the donors'.

\pn{6} The courthouse clock runs now. Nobody photographs it anymore, which is the best thing that could be said about it. The bakery still sells the cookie, though, with the hands at 4:17: a reminder, if anyone wants one, of what eleven years of neglect looked like, and of how little it cost to end.
\end{passage}

\begin{q}{3}{3.3 Structure}{The opening story becomes the small example of the essay's larger claim.}
Within the passage as a whole, paragraphs 1 and 2 function primarily as
\ch{a summary of the counterargument the writer will refute}{statistical evidence for the claim in paragraph 4}{an extended anecdote that the writer later presents as an example of a broader pattern}{a digression from the writer's main argument}
\end{q}

\begin{q}{1}{3.2 Methods of development}{The cookie shows the town had come to accept, even celebrate, the broken clock.}
The details about the bakery's cookie in paragraph 1 primarily serve to
\ch{show how thoroughly the town had come to accept, and even celebrate, the broken clock}{criticize the bakery for profiting from the county's neglect}{provide a precise date for the clock's failure}{suggest that the bakery was responsible for the clock's repair}
\end{q}

\begin{q}{2}{2.2 Evidence}{Four days, a small budget, and a quarter-sized gear stress how little the repair required.}
The details in paragraph 2 emphasize primarily
\ch{the difficulty of finding replacement parts for any machine}{how small the cost and effort of the repair were compared with the length of the neglect}{the technical skill required to repair old clocks}{the history of the company that built the clock}
\end{q}

\begin{q}{4}{3.3 Structure}{``I tell this story because\ldots'' pivots from narrative to the general argument.}
The first sentence of paragraph 3 functions mainly as
\ch{a restatement of the details in paragraph 2}{a concession that the story is unusual}{the conclusion of the essay}{a transition from the narrative to the essay's general argument}
\end{q}

\begin{q}{2}{3.2 Methods of development}{Bridge, library roof, and boiler contrast celebrated building with ignored upkeep.}
The examples in paragraph 3 (the bridge, the library roof, the boiler) are developed mainly through
\chh{classification}{contrast}{definition}{chronological narration}
\end{q}

\begin{q}{1}{3.2 Methods of development}{Deferred maintenance leads to rising repair costs: cause and effect.}
Paragraph 4 is developed chiefly through
\chh{cause and effect}{personal anecdote}{description of a place}{definition of a term}
\end{q}

\begin{q}{4}{3.1 Line of reasoning}{Paragraph 4 shows the cost; paragraph 5 identifies its cause and proposes a remedy.}
Which best describes the relationship between paragraph 4 and paragraph 5?
\ch{Paragraph 5 provides a counterexample that undermines paragraph 4.}{Paragraph 5 restates the claim of paragraph 4 with new statistics.}{Paragraph 5 returns to the anecdote that paragraph 4 interrupted.}{Paragraph 4 establishes the cost of the habit; paragraph 5 identifies its underlying cause and proposes a remedy.}
\end{q}

\begin{q}{3}{3.3 Structure}{Nobody photographing a working clock illustrates maintenance invisible in success.}
The statement ``Nobody photographs it anymore, which is the best thing that could be said about it'' (paragraph 6) most directly recalls the idea that
\ch{politicians avoid ceremonies for repairs}{the bakery's cookie is no longer popular}{maintenance is invisible precisely when it succeeds}{the clock was more interesting when it was broken}
\end{q}

\begin{q}{1}{3.3 Structure}{Story, general claim, cost, cause and remedy, return to the story.}
Which best describes the overall structure of the passage?
\ch{An anecdote, a general claim drawn from it, evidence of the claim's cost, a proposed remedy, and a return to the anecdote}{A chronological history of the courthouse from 1931 to the present}{A comparison of two clocks followed by a recommendation}{A thesis, three supporting examples, and a summary of each}
\end{q}

\section{Passage 2: Basic research (reading)}

\begin{passage}{Passage 2}{The following passage is from a lecture delivered in 2022 to a committee of state legislators by a biologist who directs a university research institute.}
\pn{1} You have asked me to explain why the state should pay for research that has no obvious use. That is a fair question, and I want to answer it honestly, which means beginning with a confession: most of that research will turn out to have no use at all.

\pn{2} Let me define my terms. By ``basic research'' I do not mean research that is easy, or research that is unimportant. I mean research whose purpose is to understand something rather than to make something. The scientist who studies how a jellyfish glows is doing basic research. The engineer who designs a better light bulb is not.

\pn{3} Here is the difficulty: the two are connected, but the connection is visible only in hindsight. In the 1960s, researchers studying a glowing jellyfish isolated the protein responsible. They wanted to know how the animal made light; no one imagined a use. Some thirty years later, other scientists realized that the protein could be attached to other proteins, making them visible under a microscope. Today it is used in laboratories around the world to watch cancer cells spread and neurons grow. No committee asked to fund a tool for tracking cancer would have thought to start with a jellyfish.

\pn{4} That is the pattern, not the exception. The first lasers were built by physicists curious about how atoms release light, years before anyone thought to scan a barcode or reshape a cornea with one.

\pn{5} I will grant that this argument can be abused. ``You never know what might be useful'' can justify anything, and you are right to demand more than that. So I would ask you to judge us not by whether each project pays off, which no one can know in advance, but by whether the work is rigorous, whether it is shared openly, and whether, across a decade, our work as a whole produces discoveries that others build on. That is how you would judge an investment fund. It is fair to judge us the same way.

\pn{6} The question, then, is not whether the state can afford to pay for research without obvious uses. It is whether a state that wants the uses can afford not to.
\end{passage}

\begin{q}{3}{2.4 Concession}{Conceding a skeptic's point up front builds credibility.}
The ``confession'' in paragraph 1 primarily serves to
\ch{suggest that the speaker opposes funding basic research}{introduce the example of the jellyfish}{build credibility by conceding at the outset a point skeptics would raise}{criticize the legislators' question as unfair}
\end{q}

\begin{q}{2}{3.2 Methods of development}{Defining the key term clarifies what she defends and heads off a misunderstanding.}
Paragraph 2 is primarily an example of
\ch{classification of the sciences into fields}{definition used to clarify what the speaker is defending and to prevent a misunderstanding}{narration that establishes the speaker's background}{comparison of two research institutes}
\end{q}

\begin{q}{3}{3.2 Methods of development}{The jellyfish story narrates how curiosity led to an unforeseen use.}
Paragraph 3 develops the speaker's argument chiefly through
\ch{an analogy between science and agriculture}{a set of statistics on research funding}{a narrative example showing how curiosity-driven research led to an unforeseen use}{a definition of cancer research}
\end{q}

\begin{q}{4}{2.3 Commentary}{The last sentence draws out the implication: targeted funding would have missed it.}
The final sentence of paragraph 3 (``No committee \ldots\ with a jellyfish'') primarily serves to
\ch{criticize the scientists who studied the jellyfish}{concede that the jellyfish research was wasteful}{suggest that cancer research is overfunded}{draw out the implication that funding only for known uses would have missed the discovery}
\end{q}

\begin{q}{4}{3.4 Transitions}{``That is the pattern'' generalizes from one example to a broader claim.}
The sentence ``That is the pattern, not the exception'' (paragraph 4) functions as
\ch{a call for the legislators to act}{a summary of the speaker's credentials}{a concession that the jellyfish case was unusual}{a transition that generalizes from one example to a broader claim}
\end{q}

\begin{q}{1}{2.4 Concession}{She grants the argument's weakness and offers a standard to answer it.}
In paragraph 5, the speaker primarily
\ch{acknowledges a weakness in her argument and proposes a standard that addresses it}{argues that every project must show results within a year}{abandons the argument that basic research is valuable}{shifts the discussion to the state's investment funds}
\end{q}

\begin{q}{1}{4.3 Comparisons}{The investment-fund standard is familiar: judge the portfolio, not each project.}
The comparison to ``an investment fund'' (paragraph 5) primarily serves to
\ch{offer lawmakers a familiar way to judge research as a whole rather than project by project}{suggest that research institutes should invest in the stock market}{warn that research funding is risky and should be reduced}{show that scientists understand finance better than legislators do}
\end{q}

\begin{q}{2}{3.3 Structure}{The ending reverses the opening question, restating the thesis.}
The final paragraph functions primarily to
\ch{concede that the state cannot afford basic research}{reverse the terms of the question posed in paragraph 1, restating the speaker's position}{summarize the jellyfish example}{introduce a new example of basic research}
\end{q}

\section{Passage 3: Carver Park Pool (writing)}

\begin{passage}{Passage 3}{The following is a draft of an essay a student is writing for a local history contest. It may need revision.}
\pn{1} \sn{1} For most of the twentieth century, Carver Park Pool was the busiest place in Millbrook from June to August. \sn{2} It opened in 1938, built by workers paid through a federal jobs program. \sn{3} The pool held eight hundred swimmers, and on hot Saturdays the line stretched past the bandstand.

\pn{2} \sn{4} In 1971, the city closed the pool, citing the cost of repairing its cracked basin. \sn{5} For the next two decades, Millbrook had no public pool at all. \sn{6} Children who could not afford the private swim club swam in the old quarry, where three drowned between 1974 and 1989.

\pn{3} \sn{7} The pool's history was not always a happy one. \sn{8} Until 1953, Black residents were permitted to swim only on Monday afternoons, just before the water was drained and replaced. \sn{9} That rule ended after a group of parents, led by a pharmacist named Lucille Ames, filled the city council chamber for six meetings in a row.

\pn{4} \sn{10} In 1994, a new pool opened on the same site. \sn{11} It was named for Lucille Ames. \sn{12} Today it hosts free swim lessons for every second grader in the district.

\pn{5} \sn{13} The story of Carver Park Pool is a story about what a town decides it owes its children.
\end{passage}

\begin{q}{3}{3.5 Writing: placement}{The 1953 events come before the 1971 closing, so paragraph 3 belongs before paragraph 2.}
To make the essay's organization chronological, where should paragraph 3 be placed?
\ch{Where it is now}{Before paragraph 1}{Before paragraph 2}{After paragraph 4}
\end{q}

\begin{q}{4}{3.3 Structure}{A topic sentence should name the paragraph's focus, unequal access, and contrast with the busy pool just described.}
Assuming paragraph 3 is placed immediately after paragraph 1, which version of sentence 7 best introduces the paragraph's focus?
\ch{(as it is now)}{Swimming was very popular in Millbrook.}{Lucille Ames was a pharmacist.}{For its first fifteen years, however, the pool was not open to everyone equally.}
\end{q}

\begin{q}{2}{3.5 Writing: placement}{A detail about roads and post offices shifts focus away from the pool.}
The writer is considering adding the following sentence after sentence 2.

\emph{The same jobs program also built roads, bridges, and post offices across the country.}

Should the writer make this addition?
\ch{Yes, because it explains how the pool was paid for.}{No, because it shifts the focus away from the history of the pool.}{Yes, because it provides a transition to paragraph 2.}{No, because it contradicts sentence 2.}
\end{q}

\begin{q}{2}{3.4 Transitions}{Swimming in the quarry follows from having no public pool: a consequence.}
Which transition, if placed at the beginning of sentence 6 (adjusting capitalization as needed), best connects it to sentence 5?
\chh{Nevertheless,}{As a result,}{In contrast,}{Similarly,}
\end{q}

\begin{q}{3}{4.2 Syntax}{Folding the name into the opening sentence is concise and keeps the focus on the reopening.}
The writer wants to combine sentences 10 and 11. Which choice most effectively combines them?
\ch{A new pool, in 1994, opened, named for Lucille Ames, on the same site.}{In 1994, a new pool opened on the same site; it was named; for Lucille Ames.}{In 1994, a new pool named for Lucille Ames opened on the same site.}{In 1994, a new pool opened on the same site, and also the new pool was named for Lucille Ames.}
\end{q}

\begin{q}{2}{3.5 Writing: placement}{Where the money went belongs right after the closing it explains.}
The writer wants to add the following sentence to the essay.

\emph{The money set aside for the pool's repair was spent instead on a parking garage downtown.}

Where would the sentence best be placed?
\chh{After sentence 1}{After sentence 4}{After sentence 8}{After sentence 12}
\end{q}

\begin{q}{4}{3.3 Structure}{A thesis belongs at the end of the introduction, after the opening scene.}
The writer wants to add the following sentence to paragraph 1 to state the essay's main claim.

\emph{Its history, from its exclusions to its closing and its rebirth, traces what the town believed it owed its children.}

Where would the sentence best be placed?
\chh{Before sentence 1}{After sentence 1}{After sentence 2}{After sentence 3}
\end{q}

\begin{q}{1}{3.3 Structure}{The conclusion should tie the closing and the reopening back to the essay's claim.}
The writer wants to add a sentence after sentence 13 that connects the events of the essay to its main claim. Which choice best accomplishes this goal?
\ch{Closing the pool and reopening it were both answers to that question, and only one of them was worth the cost.}{The new pool is twenty-five meters long and has six lanes.}{Swimming is excellent exercise for people of all ages.}{Many towns across the country have public pools.}
\end{q}

\section{Free response}

\begin{frq}{Question 1 \textperiodcentered\ Synthesis \textperiodcentered\ 6 points \textperiodcentered\ 55 minutes (15 reading + 40 writing)}
\small
\textbf{Directions.} The following prompt is based on the accompanying six sources. The sources, their authors, and the district and school named in them are fictional, written for this practice set.

Many high schools have adopted policies restricting student phone use. Some require phones to be locked away for the entire school day (``bell to bell''); others restrict phones only during class time; still others leave the decision to individual teachers.

\textbf{Prompt.} Carefully read the following six sources, including the introductory information for each source. Write an essay that synthesizes material from at least three of the sources and develops your position on the extent to which high schools should restrict student phone use during the school day.

Source A (Okafor) \quad Source B (Ridgeview survey) \quad Source C (Ferreira) \quad Source D (student editorial) \quad Source E (policy brief) \quad Source F (Westfield logs)

In your response you should do the following:
\begin{itemize}[leftmargin=1.5em,itemsep=1pt]
\item Respond to the prompt with a thesis that presents a defensible position.
\item Select and use evidence from at least three of the provided sources to support your line of reasoning. Indicate clearly the sources used through direct quotation, paraphrase, or summary. Sources may be cited as Source A, Source B, etc., or by using the description in parentheses.
\item Explain how the evidence supports your line of reasoning.
\item Use appropriate grammar and punctuation in communicating your argument.
\end{itemize}
\end{frq}

\needspace{10\baselineskip}
\begin{source}{Source A}{Okafor, Daniel. ``I Teach. Take the Phones.'' Opinion column by a high school English teacher, 2024.}
I have taught for nineteen years, and I can tell you the exact moment my classroom changed: the year every student arrived with a device designed by some of the most talented engineers on earth to pull their attention away from whatever is in front of them, including me. When phones are allowed in pockets but ``not out,'' I spend the period policing the line between pocket and desk. When phones are locked away for the whole day, something I had nearly forgotten returns: lunchrooms get loud with conversation, and students stay with a difficult paragraph long enough to understand it. A rule that applies only during class leaves the hallway and lunch as phone time, and students arrive at class still inside whatever they were scrolling.
\end{source}

\needspace{10\baselineskip}
\begin{source}{Source B}{``One Year of Bell-to-Bell.'' Survey of 1,200 students conducted by Ridgeview Unified School District, 2024, one year after the district required phones to be locked in pouches for the full school day.}
\begin{center}\small
\begin{tabular}{lccc}
\toprule
Statement & Agree & Neutral & Disagree \\
\midrule
I can focus better in class than last year & 54\% & 28\% & 18\% \\
I talk more with friends at lunch & 61\% & 24\% & 15\% \\
I feel less anxious during the school day & 38\% & 33\% & 29\% \\
I have had trouble reaching a family member when I needed to & 27\% & 19\% & 54\% \\
I would keep the policy & 41\% & 22\% & 37\% \\
\bottomrule
\end{tabular}
\end{center}
\end{source}

\needspace{10\baselineskip}
\begin{source}{Source C}{Ferreira, Maya. ``Restriction Is Not a Skill.'' Article by a pediatrician, published in a magazine for parents, 2023.}
The evidence that constant phone use harms adolescents' sleep and attention is strong enough that I support limits during class. But I worry when schools treat removal as the whole solution. A student whose phone is locked away from 8:00 to 3:00 has not learned to put it down; she has learned that someone else will put it down for her. At 3:01 she is on her own, and she will be on her own at nineteen. The schools I have seen succeed pair restrictions in class with explicit teaching about attention: why notifications are designed the way they are, how to turn them off, how to notice the pull.
\end{source}

\needspace{10\baselineskip}
\begin{source}{Source D}{``Not Every Student's Phone Is a Toy.'' Editorial by the student editorial board of a high school newspaper, 2024.}
The debate over phones assumes every student has the same day. Many of us don't. Some of us coordinate rides for younger siblings, some of us work afternoon shifts that change by text, and some of us use apps that help us read, track blood sugar, or translate for our parents. A bell-to-bell ban treats all of that as distraction. We support phones away during class; we've seen what they do to a lecture. But lunch and passing periods belong to us, and a policy that can't tell the difference between a game and a glucose monitor is not a good policy.
\end{source}

\needspace{10\baselineskip}
\begin{source}{Source E}{``Phone Policies: What Enforcement Costs.'' Policy brief from a state association of school administrators, 2024.}
Lockable pouches cost roughly \$20 to \$30 per student, plus replacements; a school of 1,500 students can expect first-year costs of \$35,000 to \$45,000. The larger cost is staff time: schools in our survey reported that teachers and administrators spent the equivalent of one to two full-time positions on enforcement during the first semester, falling by about half in the second. Policies that were applied inconsistently, enforced strictly by some teachers and ignored by others, produced the most complaints from both students and parents.
\end{source}

\needspace{10\baselineskip}
\begin{source}{Source F}{Teacher logs from Westfield High School, 2023--2024. Average minutes per 50-minute class period spent on phone-related interruptions (reminders, confiscations, off-task time) before and after a class-time-only phone policy.}
\begin{center}\small
\begin{tabular}{lcc}
\toprule
Subject & Before policy & After policy \\
\midrule
English & 7.5 & 2.0 \\
Mathematics & 6.0 & 1.5 \\
Science & 5.5 & 2.5 \\
Social studies & 8.0 & 2.0 \\
\bottomrule
\end{tabular}
\end{center}
\end{source}

\makekey

\subsection*{Essay scoring guide (6 points)}
\small
The rows follow the published AP English Language synthesis rubric (thesis 1, evidence and commentary 4, sophistication 1). Row B scores of 2--4 require evidence from at least three sources; fewer than two earns 0.
\begin{compare}{Row}{Earns credit}{Common miss}
A. Thesis (1) & A defensible position on the extent of restriction, e.g.\ ``High schools should keep phones out of classrooms entirely but leave lunch and passing periods alone, pairing the rule with lessons on attention, because the gains in focus come from class-time limits, while all-day bans add cost and burden students who rely on their phones.'' & ``Phones have pros and cons''; ``this essay will discuss phone policies''; a thesis that ignores the question of extent. \\
B. Evidence and commentary (4) & Evidence from three or more sources, grouped by idea and explained, e.g.\ Source F (interruptions fell from 6--8 to about 2 minutes) shows class-time rules recover instruction; Source B shows real social gains from bell-to-bell (61\% talk more at lunch) but only 41\% would keep it and 27\% had trouble reaching family; Source D explains why: phones serve work, family, and medical needs outside class; Source E shows the cost and staff time of all-day enforcement; Source C argues restriction must be paired with teaching self-regulation. Commentary explains what each piece of evidence shows about \emph{how much} schools should restrict. 4: all claims supported and consistently explained. 3: explanation for only some. 2: three sources with little line of reasoning. 1: two sources, summarized. & One paragraph per source, summarized in order; quoting Source B's table without interpreting what the numbers mean; using only Sources A and F. \\
C. Sophistication (1) & Weighs competing values the sources represent (attention vs.\ autonomy and equity), answers a source that cuts against the thesis (e.g.\ concedes Okafor's point that hallway scrolling carries into class, then answers it with Ferreira), notes a limitation (Source B is self-reported and from one district; Source F covers only one school), or sustains a vivid, persuasive style. & Mentioning ``both sides'' in one sentence without developing the tension. \\
\end{compare}

\end{document}
